% GENERATED FILE. Edit canonical lessons, then run npm run generate:books.
% canonical-source-hash: fnv1a64:7b64654ded5ba8f3
% canonical-lessons: GE-C400-scheinwerfer, GE-C400-windschutzscheibe, GE-C400-scheibenwischer, GE-C400-rueckspiegel, GE-C400-blinken

\chapter{Freie Sicht}
\label{ch:ge-freie-sicht}
\hlchaptermodality{400}

\begin{chapteropening}
\textbf{By the end of this chapter:} \emph{I can talk about seeing and signalling from a car: the headlights, the windscreen, the wipers, the rear-view mirror and indicating.}

Complete the last lesson of chapter 400: I can talk about seeing and signalling from a car: the headlights, the windscreen, the wipers, the rear-view mirror and indicating.
\end{chapteropening}

\section[der Scheinwerfer]{der Scheinwerfer — the headlight; the spotlight}

\label{lesson:GE-C400-scheinwerfer}



\begin{quote}
Before the new one: say the German for the brake, then the German for to change gear.
\end{quote}

\subsection*{You'll want to know: der Scheinwerfer}
\textbf{der Scheinwerfer} — \textquotedblleft{}the headlight; the spotlight\textquotedblright{}.

\begin{quote}
Am Abend schaltet sie die Scheinwerfer ein.
\end{quote}

\begin{quote}
\emph{In the evening she switches the headlights on.}
\end{quote}

Same form in the plural: \emph{die Scheinwerfer}. On a stage, the spotlight is also \emph{der Scheinwerfer}.

\begin{cousinweb}
\emph{Scheinwerfer} is built from \emph{Schein} 'light, glow' and \emph{Werfer} 'thrower', from \emph{werfen} 'to throw': a thrower of light. Both are inherited: \emph{Schein} goes back to Proto-Germanic \emph{skīnaną} 'to shine' and \emph{werfen} to \emph{werpaną} 'to throw' (EtymDB).
\end{cousinweb}

\begin{grammarlens}[title={What you've built}]
One more word for telling, explaining and describing at B1.
\end{grammarlens}

\subsection*{Your turn}
\begin{itemize}
\raggedright
  \item \emph{Say it:} \emph{der Scheinwerfer}
  \item \emph{Say it:} \emph{der Scheinwerfer}, once more
  \item \emph{Say it:} a sentence of your own with \emph{der Scheinwerfer}
  \item \emph{Recall:} say the German for the brake, then the German for to change gear, then say \emph{der Scheinwerfer} again
\end{itemize}

\begin{tcolorbox}[breakable,colback=teal!4,colframe=teal!35!black,title={Before you move on}]
What does der Scheinwerfer mean? (\textquotedblleft{}The headlight; the spotlight\textquotedblright{}.) Say it once more.
\end{tcolorbox}
\section[die Windschutzscheibe]{die Windschutzscheibe — the windscreen (US: windshield)}

\label{lesson:GE-C400-windschutzscheibe}



\begin{quote}
Before the new one: say the German for to change gear, then the German for the headlight.
\end{quote}

\subsection*{You'll want to know: die Windschutzscheibe}
\textbf{die Windschutzscheibe} — \textquotedblleft{}the windscreen (US: windshield)\textquotedblright{}.

\begin{quote}
Ein kleiner Stein hat die Windschutzscheibe getroffen.
\end{quote}

\begin{quote}
\emph{A small stone hit the windscreen.}
\end{quote}

Plural: \emph{die Windschutzscheiben}. Three parts: \emph{Wind} + \emph{Schutz} + \emph{Scheibe}.

\begin{cousinweb}
\emph{Windschutzscheibe} is built from \emph{Wind}, \emph{Schutz} 'protection' and \emph{Scheibe} 'pane, disc', all inherited Germanic words: \emph{Wind} goes back to Proto-Germanic \emph{windaz} and \emph{Scheibe} to Old High German \emph{scība} (EtymDB).
\end{cousinweb}

\begin{grammarlens}[title={What you've built}]
One more word for telling, explaining and describing at B1.
\end{grammarlens}

\subsection*{Your turn}
\begin{itemize}
\raggedright
  \item \emph{Say it:} \emph{die Windschutzscheibe}
  \item \emph{Say it:} \emph{die Windschutzscheibe}, once more
  \item \emph{Say it:} a sentence of your own with \emph{die Windschutzscheibe}
  \item \emph{Recall:} say the German for to change gear, then the German for the headlight, then say \emph{die Windschutzscheibe} again
\end{itemize}

\begin{tcolorbox}[breakable,colback=teal!4,colframe=teal!35!black,title={Before you move on}]
What does die Windschutzscheibe mean? (\textquotedblleft{}The windscreen (US: windshield)\textquotedblright{}.) Say it once more.
\end{tcolorbox}
\section[der Scheibenwischer]{der Scheibenwischer — the windscreen wiper}

\label{lesson:GE-C400-scheibenwischer}



\begin{quote}
Before the new one: say the German for the headlight, then the German for the windscreen.
\end{quote}

\subsection*{You'll want to know: der Scheibenwischer}
\textbf{der Scheibenwischer} — \textquotedblleft{}the windscreen wiper\textquotedblright{}.

\begin{quote}
Bei starkem Regen laufen die Scheibenwischer sehr schnell.
\end{quote}

\begin{quote}
\emph{In heavy rain the wipers go very fast.}
\end{quote}

Same form in the plural. Often shortened to \emph{der Wischer}.

\begin{cousinweb}
\emph{Scheibenwischer} is a compound of \emph{Scheibe} 'pane' and \emph{Wischer} 'wiper', from \emph{wischen} 'to wipe', which is inherited through Middle High German \emph{wischen} from Old High German \emph{wisken} (EtymDB).
\end{cousinweb}

\begin{grammarlens}[title={What you've built}]
One more word for telling, explaining and describing at B1.
\end{grammarlens}

\subsection*{Your turn}
\begin{itemize}
\raggedright
  \item \emph{Say it:} \emph{der Scheibenwischer}
  \item \emph{Say it:} \emph{der Scheibenwischer}, once more
  \item \emph{Say it:} a sentence of your own with \emph{der Scheibenwischer}
  \item \emph{Recall:} say the German for the headlight, then the German for the windscreen, then say \emph{der Scheibenwischer} again
\end{itemize}

\begin{tcolorbox}[breakable,colback=teal!4,colframe=teal!35!black,title={Before you move on}]
What does der Scheibenwischer mean? (\textquotedblleft{}The windscreen wiper\textquotedblright{}.) Say it once more.
\end{tcolorbox}
\section[der Rückspiegel]{der Rückspiegel — the rear-view mirror}

\label{lesson:GE-C400-rueckspiegel}



\begin{quote}
Before the new one: say the German for the windscreen, then the German for the windscreen wiper.
\end{quote}

\subsection*{You'll want to know: der Rückspiegel}
\textbf{der Rückspiegel} — \textquotedblleft{}the rear-view mirror\textquotedblright{}.

\begin{quote}
Vor dem Abbiegen schaut die Fahrerin in den Rückspiegel.
\end{quote}

\begin{quote}
\emph{Before turning off, the driver looks in the rear-view mirror.}
\end{quote}

Same form in the plural. The mirrors on the sides are \emph{die Außenspiegel}.

\begin{cousinweb}
\emph{Rückspiegel} is built from \emph{rück-} 'back' and \emph{Spiegel} 'mirror'. \emph{Spiegel} is an early borrowing from Latin \emph{speculum} 'mirror', built on \emph{specere} 'to look' (Lewis \& Short); EtymDB lists Middle High German \emph{spiegel}.
\end{cousinweb}

\begin{grammarlens}[title={What you've built}]
One more word for telling, explaining and describing at B1.
\end{grammarlens}

\subsection*{Your turn}
\begin{itemize}
\raggedright
  \item \emph{Say it:} \emph{der Rückspiegel}
  \item \emph{Say it:} \emph{der Rückspiegel}, once more
  \item \emph{Say it:} a sentence of your own with \emph{der Rückspiegel}
  \item \emph{Recall:} say the German for the windscreen, then the German for the windscreen wiper, then say \emph{der Rückspiegel} again
\end{itemize}

\begin{tcolorbox}[breakable,colback=teal!4,colframe=teal!35!black,title={Before you move on}]
What does der Rückspiegel mean? (\textquotedblleft{}The rear-view mirror\textquotedblright{}.) Say it once more.
\end{tcolorbox}
\section[blinken]{blinken — to indicate (in a car); to flash, to twinkle}

\label{lesson:GE-C400-blinken}



\begin{quote}
Before the new one: say the German for the windscreen wiper, then the German for the rear-view mirror.
\end{quote}

\subsection*{You'll want to know: blinken}
\textbf{blinken} — \textquotedblleft{}to indicate (in a car); to flash, to twinkle\textquotedblright{}.

\begin{quote}
Er blinkt rechts und fährt in die kleine Straße.
\end{quote}

\begin{quote}
\emph{He indicates right and turns into the little street.}
\end{quote}

\emph{Er blinkt}, \emph{er hat geblinkt}. Stars also \emph{blinken}: they twinkle.

\begin{cousinweb}
\emph{blinken} is inherited: EtymDB traces it through Middle Low German \emph{blinken} to Proto-Germanic \emph{blinkaną} 'to shine, to gleam'. English \emph{blink} is a cousin.
\end{cousinweb}

\begin{grammarlens}[title={What you've built}]
One more word for telling, explaining and describing at B1.
\end{grammarlens}

\subsection*{Your turn}
\begin{itemize}
\raggedright
  \item \emph{Say it:} \emph{blinken}
  \item \emph{Say it:} \emph{blinken}, once more
  \item \emph{Say it:} a sentence of your own with \emph{blinken}
  \item \emph{Recall:} say the German for the windscreen wiper, then the German for the rear-view mirror, then say \emph{blinken} again
\end{itemize}

\begin{tcolorbox}[breakable,colback=teal!4,colframe=teal!35!black,title={Before you move on}]
What does blinken mean? (\textquotedblleft{}To indicate (in a car); to flash, to twinkle\textquotedblright{}.) Say it once more.
\end{tcolorbox}
